\subsubsection{Interfaces Externas - El vínculo con usuarios, academia y HPC [Pendientes: 1. Confirmar participación de John Jairo Corredor Franco. 2. Identificar usuarios evaluadores. 3. Precisar los datos de contacto y medios de comunicación de todas las personas y entidades participantes.]}\label{sec:spmp-8-4-1}
La presente sección representa aquellas personas y entidades que no participan en el desarrollo del proyecto al mismo nivel o de forma tan directa como lo hacen los cuatro estudiantes pero que aportan en el mismo y son indispensables para su correcta conclusión.
\begin{longtable}{|
p{0.20\dimexpr\linewidth-8\tabcolsep-5\arrayrulewidth\relax}|
p{0.25\dimexpr\linewidth-8\tabcolsep-5\arrayrulewidth\relax}|
p{0.35\dimexpr\linewidth-8\tabcolsep-5\arrayrulewidth\relax}|
p{0.20\dimexpr\linewidth-8\tabcolsep-5\arrayrulewidth\relax}|}
\caption{Participantes externos y ámbitos de relación. Fuente: identificación de participantes del proyecto.}\\\hline
  \textbf{Persona o entidad} & \textbf{Descripción} & \textbf{Responsabilidades con el proyecto} & \textbf{Canal de comunicación}\\\hline
\endfirsthead\hline
\textbf{Persona o entidad} & \textbf{Descripción} &\textbf{Responsabilidades con el proyecto} & \textbf{Canal de comunicación}\\\hline
\endhead
Jaime Andrés Pavlich Mariscal & Profesor de la Pontificia Universidad Javeriana, Director del proyecto de grado y acompañamiento académico.& Revisar el avance al cierre de cada iteración, orientar las decisiones del proyecto y participar en la aceptación del producto (8.3) & Correo institucional y reuniones pre-programadas\\\hline
Juan Guillermo Torres Hurtado & Coordinación de HPC-ZINE. & Capacitación (ya hecha), cotización, acceso individual en 2027-1, asignación de nodos por Slurm y soporte para la instalación de FDS, que hoy no está instalado & Correo institucional y reuniones pre-programadas\\\hline
Andrés Oswaldo Calderón Romero & Profesor de la Pontificia Universidad Javeriana, asesor del proyecto & Dar acompañamiento técnico, en las areas de computación distribuida y manejo de datos geograficos & Correo institucional y reuniones pre-programadas\\\hline
Bomberos de Bogotá & Usuarios que representan al comandante o coordinador de la emergencia.& Evaluar la versión final en Semi-transición (I7, del 26 de abril al 9 de mayo) y participar en la aceptación (8.3) & WhatsApp, llamadas telefónicas y reuniones pre-programadas\\\hline
  Coordinación de Trabajo de Grado de la Carrera de Ingeniería de Sistemas & Instancia académica de la Pontificia Universidad Javeriana que gestiona los trabajos de grado. & Establecer las fechas de entrega y sustentación, recibir la entrega final en la semana 17 del periodo 2027-1 y gestionar la evaluación del trabajo de grado & Correo institucional\\\hline
IDECA y Jardín Botánico de Bogotá & Proveedores de datos geográficos abiertos de Bogotá (cobertura de la tierra, terreno y cartografía de referencia). & Sin responsabilidades directas con el proyecto; el proyecto depende de la disponibilidad de sus datos para preparar las entradas de FDS & Portales web y servicios de descarga de datos abiertos\\\hline
Open-Meteo & Proveedor de datos meteorológicos abiertos mediante API. & Sin responsabilidades directas con el proyecto; el proyecto depende de la disponibilidad de su API para obtener las condiciones meteorológicas de los escenarios y de la consulta en la aplicación & Portal web y documentación de la API\\\hline
\end{longtable}
